% !TeX root = semana05.tex
% !TeX program = pdflatex
% Dos clases de 2 horas: 1 hora de corrección y 3 horas de continuidad.
\documentclass[10pt,a5paper]{article}
\newif\ifmostrarSoluciones
\mostrarSolucionestrue
\usepackage{tikz}
\usepackage{estilo_apuntes}
\usetikzlibrary{arrows.meta,positioning,babel}
\encabezado{Semana 5}{Apunte semanal}

\begin{document}
\tituloclase{Continuidad de funciones}{Continuidad puntual, discontinuidades y operaciones}{Semana 5 · 4 horas}

\begin{objetivos}
\begin{itemize}
  \item Distinguir el valor de una función, su límite y su continuidad en un punto.
  \item Estudiar continuidad en intervalos, considerando los extremos incluidos.
  \item Clasificar discontinuidades y determinar cuáles admiten una extensión continua.
  \item Ajustar parámetros en funciones por tramos mediante límites laterales.
  \item Justificar continuidad mediante operaciones y composición, con sus restricciones de dominio.
  \item Calcular límites por continuidad y cambio de variable, verificando las hipótesis.
\end{itemize}
\end{objetivos}

\begin{resumen}
Una función es continua en un punto cuando su valor coincide con el límite allí. En funciones por tramos se comparan los límites laterales con el valor asignado. Las reglas de continuidad permiten después estudiar expresiones más complejas a partir de funciones conocidas, siempre dentro de sus dominios.
\end{resumen}


\section*{Continuidad en un punto}

\begin{definicion}[continuidad en un punto interior]
Sea $a$ un punto interior del dominio de $f$. La función $f$ es \textbf{continua en $a$} si
\[
  \lim_{x\to a}f(x)=f(a).
\]
Esto exige verificar tres condiciones:
\begin{enumerate}
  \item $f(a)$ está definido.
  \item Existe el límite bilateral $\lim_{x\to a}f(x)$ como número real.
  \item Ese límite es igual a $f(a)$.
\end{enumerate}
\end{definicion}

Es decir, 
\[
  \lim_{x\to a^-}f(x)=\lim_{x\to a^+}f(x)=f(a).
\]
Un límite estudia puntos cercanos; la continuidad relaciona ese comportamiento con el valor en el punto.

\begin{ejemplo}[el límite existe, pero el valor no coincide]
Estudie la continuidad en $x=1$ de
\[
  h(x)=\begin{cases}
    \dfrac{x^2-1}{x-1},&x\neq1,\\[1mm]
    5,&x=1.
  \end{cases}
\]
¿Qué valor se debería asignar a $h(1)$ para obtener continuidad?
\end{ejemplo}
\ifmostrarSoluciones
\begin{solucion}
Para $x\neq1$, $h(x)=x+1$. Por tanto,
\[
  \lim_{x\to1}h(x)=2\neq5=h(1).
\]
El límite existe y $h(1)$ está definido, pero falla la tercera condición. La función es discontinua en $1$. Al reemplazar únicamente $h(1)$ por $2$, se obtiene una función continua en ese punto.
\end{solucion}
\else\vspace{18mm}\fi

\begin{center}
\begin{tikzpicture}[x=1.2cm,y=0.62cm,every node/.style={font=\scriptsize}]
  \draw[->,Gris] (-1.4,0)--(3.1,0) node[right] {$x$};
  \draw[->,Gris] (0,-0.4)--(0,5.8) node[above] {$y$};
  \draw[Azul,thick] (-1.2,-0.2)--(2.8,3.8);
  \draw[Gris,dashed] (1,0)--(1,5);
  \filldraw[fill=white,draw=Azul,thick] (1,2) circle (2.5pt);
  \fill[Rojo] (1,5) circle (2.5pt);
  \node[below] at (1,0) {$1$};
  \node[Azul,right] at (1.15,2) {límite: $2$};
  \node[Rojo,right] at (1.15,5) {$h(1)=5$};
\end{tikzpicture}
\end{center}

\subsection*{Clasificación de discontinuidades}

Para clasificar el comportamiento cerca de $a$, primero se calculan los límites laterales. La clasificación no depende solamente de si una fórmula tiene denominador cero.

\begin{definicion}[tipos de discontinuidad]
\begin{itemize}
  \item \textbf{Removible o evitable:} existe un límite bilateral finito $L$, pero el valor en $a$ no está definido o es distinto de $L$. Se corrige asignando $f(a)=L$.
  \item \textbf{De salto:} los dos límites laterales existen y son finitos, pero son distintos. Cambiar solamente $f(a)$ no elimina el salto.
  \item \textbf{Infinita:} al menos uno de los límites laterales es $+\infty$ o $-\infty$. No se corrige asignando un valor real en $a$.
  \item \textbf{Oscilatoria:} la oscilación impide que exista alguno de los límites laterales, sin que ese comportamiento sea un límite infinito.
\end{itemize}
\end{definicion}

\begin{ejemplo}[]
Clasifique el comportamiento en $0$ de las funciones
\[
  s(x)=\begin{cases}0,&x<0,\\1,&x\geq0,\end{cases}
  \qquad r(x)=\frac{1}{x^2},
  \qquad o(x)=\sin(1/x).
\]
Las dos últimas están definidas para $x\neq0$.
\end{ejemplo}
\ifmostrarSoluciones
\begin{solucion}
En $s$, los límites laterales son $0$ y $1$: hay un salto. En $r$, ambos límites laterales son $+\infty$: la discontinuidad es infinita. En $o$, los valores oscilan entre $-1$ y $1$ sin aproximarse a un único número: la discontinuidad es oscilatoria. Ninguna se elimina cambiando o asignando solamente el valor en $0$.
\end{solucion}
\else\vspace{20mm}\fi

\subsection*{Ajuste de parámetros en funciones por tramos}

\begin{criterio}[continuidad en un punto de unión]
Identifique qué tramo define $f(a)$ y calcule los dos límites laterales usando los tramos que corresponden a cada lado. Luego resuelva
\[
  \lim_{x\to a^-}f(x)=f(a),\qquad
  \lim_{x\to a^+}f(x)=f(a).
\]
Si hay varios puntos de unión, las condiciones deben cumplirse simultáneamente en todos ellos.
\end{criterio}

\begin{ejemplo}[dos parámetros y un punto de unión]
Determine $a,b\in\mathbb{R}$ para que $F$ sea continua en $x=2$:
\[
  F(x)=\begin{cases}
    \dfrac{x^2-4}{x-2},&x<2,\\[1mm]
    b,&x=2,\\
    ax+1,&x>2.
  \end{cases}
\]
\end{ejemplo}
\ifmostrarSoluciones
\begin{solucion}
Para $x<2$, se puede simplificar $F(x)=x+2$. Así,
\[
  \lim_{x\to2^-}F(x)=4,\qquad
  \lim_{x\to2^+}F(x)=2a+1,\qquad F(2)=b.
\]
La continuidad exige $4=b=2a+1$, de donde
\[
  a=\frac32,\qquad b=4.
\]
Con estos valores coinciden los límites laterales y el valor de la función.
\end{solucion}
\else\vspace{24mm}\fi

\begin{ejemplo}[un parámetro que no permite continuidad]
¿Existe $c\in\mathbb{R}$ para que la siguiente función sea continua en $0$?
\[
  Q_c(x)=\begin{cases}x+1,&x<0,\\c,&x=0,\\x+2,&x>0.\end{cases}
\]
\end{ejemplo}
\ifmostrarSoluciones
\begin{solucion}
Los límites laterales son $1$ y $2$, independientemente de $c$. No existe un límite bilateral, por lo que ningún valor de $c$ produce continuidad. Igualar solamente uno de los límites a $Q_c(0)$ no basta.
\end{solucion}
\else\vspace{14mm}\fi

\section*{Intervalos y reglas de continuidad}
\subsection*{Continuidad lateral y en intervalos}

\begin{definicion}[continuidad lateral]
Una función es \textbf{continua por la derecha en $a$} si
\[
  \lim_{x\to a^+}f(x)=f(a).
\]
Es \textbf{continua por la izquierda en $b$} si
\[
  \lim_{x\to b^-}f(x)=f(b).
\]
En ambos casos el valor en el extremo debe estar definido.
\end{definicion}

\begin{definicion}[continuidad en un intervalo]
Una función es continua en un intervalo abierto si es continua en cada punto del intervalo. Es continua en $[a,b]$ si es continua en $(a,b)$, continua por la derecha en $a$ y continua por la izquierda en $b$. En un intervalo semiabierto se exige la continuidad lateral correspondiente solamente en el extremo incluido.
\end{definicion}

\begin{ejemplo}[un extremo del dominio]
La función $f(x)=\sqrt{x+2}$ tiene dominio $[-2,\infty)$. Estudie su continuidad, incluyendo el extremo $-2$.
\end{ejemplo}
\ifmostrarSoluciones
\begin{solucion}
En $(-2,\infty)$ la raíz es continua. En el extremo,
\[
  \lim_{x\to-2^+}\sqrt{x+2}=0=f(-2).
\]
Por tanto, $f$ es continua en $[-2,\infty)$. No se exige un límite desde la izquierda de $-2$, pues allí no hay puntos de su dominio.
\end{solucion}
\else\vspace{16mm}\fi

\subsection*{Álgebra de funciones continuas}

Las funciones constante $f(x)=c$ e identidad $I(x)=x$ son continuas en $\mathbb{R}$, pues sus límites en $a$ son $c$ y $a$, respectivamente. A partir de ellas se construyen los polinomios mediante las siguientes operaciones.

\begin{teorema}[operaciones que conservan continuidad]
Si $f$ y $g$ son continuas en $a$ y sea $c\in\mathbb{R}$, entonces son continuas en $a$ las funciones
\[
  f+g,\quad f-g,\quad cf,\quad fg.
\]
También $f/g$ es continua en $a$ si $g(a)\neq0$. Las potencias enteras positivas de una función continua son continuas.
\end{teorema}

Estas propiedades se obtienen del álgebra de límites. Por ejemplo, si $g(a)\neq0$,
\[
  \lim_{x\to a}\frac{f(x)}{g(x)}
  =\frac{\lim_{x\to a}f(x)}{\lim_{x\to a}g(x)}
  =\frac{f(a)}{g(a)}.
\]

\begin{proposicion}[polinomios, racionales y radicales]
\begin{enumerate}
  \item Todo polinomio es continuo en $\mathbb{R}$.
  \item Una función racional $p(x)/q(x)$ es continua en cada punto donde $q(x)\neq0$.
  \item La función $\sqrt[n]{x}$ es continua en $\mathbb{R}$ si $n$ es impar, y en $[0,\infty)$ si $n$ es par. En el extremo $0$ de una raíz par se entiende continuidad por la derecha.
\end{enumerate}
\end{proposicion}

\begin{ejemplo}[dominio e intervalos de continuidad]
Determine dónde es continua cada función:
\[
  P(x)=x^3-2x+4,\quad
  R(x)=\frac{x^2-1}{x-1},\quad
  H(x)=\frac{1}{\sqrt{4-x^2}}.
\]
\end{ejemplo}
\ifmostrarSoluciones
\begin{solucion}
$P$ es continua en $\mathbb{R}$. La función $R$ es continua en $(-\infty,1)$ y $(1,\infty)$; la igualdad $R(x)=x+1$ vale solamente para $x\neq1$. En $1$ admite una extensión continua con valor $2$.

Para $H$, se exige $4-x^2>0$, equivalente a $-2<x<2$. La raíz es continua y el denominador no se anula en ese intervalo; por ello $H$ es continua en $(-2,2)$. Los extremos no pertenecen a su dominio.
\end{solucion}
\else\vspace{22mm}\fi

\subsection*{Composición de funciones continuas}

\begin{teorema}[continuidad de la composición]
Si $f$ es continua en $a$ y $g$ es continua en $f(a)$, entonces $g\circ f$ es continua en $a$, siempre que la composición esté definida en los puntos del dominio que se consideran. En particular,
\[
  \lim_{x\to a}g(f(x))=g(f(a)).
\]
En los extremos de un dominio, los límites se toman desde los lados permitidos.
\end{teorema}

Antes de aplicar el teorema se determina
\[
  \operatorname{Dom}(g\circ f)
  =\{x\in\operatorname{Dom}(f):f(x)\in\operatorname{Dom}(g)\}.
\]

\begin{ejemplo}[una composición que exige resolver una desigualdad]
Determine el dominio y la continuidad de
\[
  K(x)=\sqrt{\frac{x-1}{x+2}}.
\]
\end{ejemplo}
\ifmostrarSoluciones
\begin{solucion}
La función interior es racional y la exterior es una raíz cuadrada. Debemos exigir
\[
  x\neq-2,\qquad \frac{x-1}{x+2}\geq0.
\]
El cociente es positivo en $(-\infty,-2)$ y $(1,\infty)$, negativo en $(-2,1)$ y cero en $1$. Así,
\[
  \operatorname{Dom}(K)=(-\infty,-2)\cup[1,\infty).
\]
En los interiores de esos intervalos, $K$ es continua por composición. En $1$,
\[
  \lim_{x\to1^+}K(x)=0=K(1).
\]
Por tanto, es continua en ambos intervalos indicados, con continuidad por la derecha en $1$. El punto $-2$ permanece excluido.
\end{solucion}
\else\vspace{26mm}\fi

\subsection*{Continuidad y cambio de variable en límites}

\begin{teorema}[paso del límite por una función continua]
Suponga que $u(x)\to L$ cuando $x\to a$ y que $\varphi$ es continua en $L$. Si $u(x)$ pertenece al dominio de $\varphi$ para los puntos cercanos considerados, entonces
\[
  \lim_{x\to a}\varphi(u(x))=\varphi(L).
\]
La función interior $u$ no necesita estar definida en $a$: basta con que tenga el límite indicado.
\end{teorema}

\begin{ejemplo}[cambio de variable con una raíz continua]
Calcule $\displaystyle\lim_{x\to2}\sqrt{x^2+5}$ e indique las funciones interior y exterior.
\end{ejemplo}
\ifmostrarSoluciones
\begin{solucion}
Con $t=x^2+5$ se tiene $t\to9$ cuando $x\to2$. La función exterior $\varphi(t)=\sqrt{t}$ es continua en $9$, de modo que
\[
  \lim_{x\to2}\sqrt{x^2+5}=\sqrt{9}=3.
\]
El radicando es positivo para todo $x$, así que no hay restricciones adicionales.
\end{solucion}
\else\vspace{16mm}\fi

\section*{Práctica y trabajo autónomo}

\begin{ejercicios}[continuidad puntual y funciones por tramos]
\begin{enumerate}
  \item Para $h_k(x)=(x^2-9)/(x-3)$ si $x\neq3$, y $h_k(3)=k$, determine $k$ para que $h_k$ sea continua en $3$.
  \item Clasifique en $0$ las discontinuidades de $|x|/x$ y $1/x^2$, ambas definidas para $x\neq0$. ¿Admiten una extensión continua en $0$?
  \item Determine $a,b$ para que la función sea continua en $\mathbb{R}$:
  \[
    f(x)=\begin{cases}
      x+2,&x<0,\\
      ax+b,&0\leq x\leq2,\\
      x^2,&x>2.
    \end{cases}
  \]
  \item Determine $c$ para que $q(x)=\sqrt{x+4}$ si $x>-4$, y $q(-4)=c$, sea continua en $[-4,\infty)$.
\end{enumerate}
\end{ejercicios}

\ifmostrarSoluciones
\begin{solucion}[de continuidad puntual]
\begin{enumerate}
  \item Para $x\neq3$, $h_k(x)=x+3$, de modo que el límite es $6$. Se requiere $k=6$.
  \item En $|x|/x$, los límites laterales son $-1$ y $1$: hay un salto. En $1/x^2$, ambos límites laterales son $+\infty$: la discontinuidad es infinita. Ninguna admite extensión continua en $0$.
  \item En $0$, la continuidad exige $b=2$. En $2$, exige $2a+b=4$. Por tanto, $a=1$, $b=2$. Cada tramo es polinómico, de modo que basta comprobar estos dos puntos de unión.
  \item Se requiere $c=0$, pues el límite por la derecha en $-4$ es $0$. En el resto del dominio, la función radical es continua.
\end{enumerate}
\end{solucion}
\fi

\begin{ejercicios}[operaciones, composición y límites]
\begin{enumerate}
  \item Determine los intervalos de continuidad de
  \[
    r(x)=\frac{x+1}{x^2-4},\qquad
    s(x)=\sqrt{9-x^2},\qquad
    v(x)=\frac{1}{\sqrt{x-1}}.
  \]
  \item Determine el dominio y la continuidad de $H(x)=\sqrt{(x+2)/(x-3)}$.
  \item Calcule por continuidad $\displaystyle\lim_{x\to1}(x^2+\sqrt{x+3})/(x+2)$ y justifique el uso del cociente.
  \item Calcule $\displaystyle\lim_{x\to0}(\sqrt{4+x^2}-2)/x^2$ mediante el cambio $t=4+x^2$. Indique el lado de aproximación de $t$.
  \item Explique el error: «$p(x)=(x^2-4)/(x-2)$ es continua en $2$ porque se simplifica a $x+2$».
\end{enumerate}
\end{ejercicios}

\ifmostrarSoluciones
\begin{solucion}[de operaciones y composición]
\begin{enumerate}
  \item $r$ es continua en $(-\infty,-2)$, $(-2,2)$ y $(2,\infty)$. La función $s$ es continua en $[-3,3]$, por la derecha en $-3$ y por la izquierda en $3$. La función $v$ exige $x-1>0$ y es continua en $(1,\infty)$.
  \item La condición $(x+2)/(x-3)\geq0$, con $x\neq3$, da $(-\infty,-2]\cup(3,\infty)$. La composición es continua en esos intervalos; en $-2$ se usa continuidad por la izquierda y $H(-2)=0$.
  \item Las funciones del numerador son continuas en $1$ y el denominador vale $3\neq0$. Por tanto, el límite es $(1+2)/3=1$.
  \item Se tiene $t\to4^+$ y, para $t>4$,
  \[
    \frac{\sqrt{t}-2}{t-4}=\frac{1}{\sqrt{t}+2}.
  \]
  Por continuidad de la expresión simplificada, el límite es $1/4$.
  \item La simplificación vale solamente para $x\neq2$; la función original no está definida en $2$. Tiene límite $4$ y admite una extensión continua asignando el valor $4$ en ese punto.
\end{enumerate}
\end{solucion}
\fi

\fuentes{\textbf{Para ampliar:} \emph{Cálculo en una variable: notas de clase 2018-B}, Resumen No.~3, pp.~4--7 (continuidad y cambio de variable); \emph{Hojas de ejercicios 2018-B}, Hoja No.~03, ejercicios 3--4 (continuidad puntual y parámetros); J. Stewart, \emph{Cálculo de una variable: trascendentes tempranas}, §2.5; G. Strang--E. Herman, \emph{Cálculo, volumen 1}, OpenStax, §2.4. Ejemplos y ejercicios de este apunte elaborados para estas sesiones.}

\end{document}
